\section{Real structures and their fixed three-spheres}\hypertarget{tag-02P9}{}\label{mapsI:real-forms}
\hypertarget{tag-01PJ}{}\label{fam:section:B.4}

\begin{definition}[Real-period data]\hypertarget{tag-IP04}{}\label{I:def:real-period-data}
Use the lattice, period matrix and clockwise root generators of
Definition~\ref{def:common-conventions}, and the affine vectors of
Definition~\ref{I:def:family-parameter}. In this section write
$d=\widehat\delta$, $g=\widehat\gamma$, $u=\widehat u$ and
$w_\Lambda=\widehat w$. An unadorned $w$ in a lattice expression denotes
$w_\Lambda$; the scalar parameter remains $w=e^{-\pi ic}$. Put
\[
 e=\Pi_c d=(0,1)^{\mathsf t},\qquad E=d\otimes\gamma,
 \qquad M_0=(A_1A_2)^{-1},\qquad D_c=D_{\rm part}+\Im c<0.
\]
The circle $\rho(s)$ is translation by $se$, $s\in\mathbb R/\mathbb Z$.
In cusp coordinates $q=1/t$ and $s=(2\pi i)^{-1}\log q$, define the
bounded matrix $C_c(q)$ by $Z_c=sB_0+C_c(q)$, where
$B_0=\left(\begin{smallmatrix}0&1\\-1&0\end{smallmatrix}\right)$.
Thus $\Pi_cg=(6\mu,-s+C_{c,21})^{\mathsf t}$.
The cusp double curves $D_1,D_2,D_{12}$ belong to the edge directions
$e_1,e_2,e_1-e_2$. The action $\rho$ fixes $D_2$, while the literal
character has a limit along $D_1$, as located by the explicit toric
coordinates in Proposition~\ref{mapsI:src:A.1.2}.
\end{definition}

\begin{proposition}[Conjugate members and antiholomorphic automorphisms]\hypertarget{tag-01QG}{}\label{fam:src:B.4.2}\label{fam:st:B.4.2}
\StatementRecord{fam:src:B.4.2}{proved}{01QG}{\SeriesRef{fam:src:B.4.1}, \SeriesRef{fam:src:A.3.1}, \SeriesRef{fam:src:B.2.1}, \SeriesRef{fam:src:A.5.1}}

Use the normalization $(\beta^{\rm part}+\tau)(p_0)=0$. There is an antiholomorphic diffeomorphism $X_c\to X_{c'}$ exactly when
\[
 c'-(1-\bar c)\in2\mathbb Z,
 \qquad\text{equivalently }w'=-\bar w.
\]
Consequently $\overline{X_c}\simeq X_{1-\bar c}$. The members admitting antiholomorphic self-maps are exactly those with $\operatorname{Re}c\in\tfrac12+\mathbb Z$, equivalently $w=-\bar w$. On each such member the group of holomorphic and antiholomorphic diffeomorphisms is
\[
 \operatorname{Aut}^{\pm}(X_c)\simeq\mathbb C^*\rtimes\mathbb Z/2,
 \qquad z\longmapsto1/\bar z
\]
for the action of the second factor. Its antiholomorphic maps form the nonidentity coset. For all other members the full group is the holomorphic automorphism group of Theorem~\ref{fam:src:B.4.1}.
\end{proposition}


\begin{proof}
\begin{enumerate}[label=\arabic*.,leftmargin=*]
\item We first construct the comparison before returning its affine data to the canonical choice. Lift conjugation of the real orbifold base to the antiholomorphic involution $\iota$ of its universal cover whose cusp branch has $\tau(\iota z)=-\overline{\tau(z)}$. A lift exists by the universal covering property; it extends at the finite orbifold points, whose local equations are $s_j^{m_j}=t_j$. Choose the usual paths from a real point $t>1$, so reflection fixes the second centre and sends the first chosen lift to $g_2^{-1}z_1$. The induced involution of the deck group is
\[
 \phi(g_1)=g_2^{-1}g_1^{-1}g_2,\qquad
 \phi(g_2)=g_2^{-1},\qquad \phi(g_0)=g_0^{-1}.
\]
These identities also follow by reflecting the lifted peripheral paths; their product and their squares agree with $g_1g_2g_0=1$. The cusp branch fixes the lift uniquely.

\item Denote the affine period transformations by $G_1,G_2$ on the triple $(\tau,\mu,\beta)$ and put $K(\tau,\mu,\beta)=(-\bar\tau,1-\bar\mu,-\bar\beta)$. Direct rational substitution gives
\[
 KG_1K=G_2^{-1}G_1^{-1}G_2,\qquad KG_2K=G_2^{-1}.
\]
For example $G_2^{-1}(\tau,\mu,\beta)=(-1/\tau,(1-\mu)/\tau,\beta+3-6(1-\mu)^2/\tau)$ and
$G_1^{-1}(\tau,\mu,\beta)=(1/(1-\tau),1-\mu/(1-\tau),\beta-2+6\mu^2/(1-\tau))$. These formulas verify every constant in the two equalities. The reflected $\mu$ is cusp bounded and has the same affine laws, so their difference is a global section of $\mathcal O_B(-p_0)$ and vanishes. After this equality, the difference between the reflected and original $\beta^{\rm part}$ is a holomorphic invariant base function, hence constant; its cusp value is zero under the stated normalization. Therefore
\[
 \mu(\iota z)=1-\bar\mu(z),\qquad
 \beta^{\rm part}(\iota z)=-\overline{\beta^{\rm part}(z)}.
\]
The original $\tau$ identity follows from its $j$ equation and its specified cusp branch. In particular $D_{-\bar c}(\iota z)=D_c(z)$, so this comparison stays in the exact admissible domain.

\item Let
\[
 H=\begin{pmatrix}1&0&0&0\\0&1&0&0\\-6&0&-1&0\\0&-1&0&-1\end{pmatrix}.
 \label{I:eq:real-reflection-matrix}
\]
Then $H^2=I$, $HA_2H=A_2^{-1}$, $HA_1H=A_2^{-1}A_1^{-1}A_2$, and
\[
 \Pi_{-\bar c}(\iota z)H=-\overline{\Pi_c(z)}.
\]
Thus $(z,\zeta)\mapsto(\iota z,-\bar\zeta)$ is a well-defined antiholomorphic torus map. It also descends under the middle deck action. Indeed covariance and the displayed lattice conjugacies imply $R_{\phi(g)}(\iota z)=\overline{R_g(z)}$ by multiplying the full real-invertible period matrices. Thus the reflection intertwines the middle deck actions.

\item At the first finite centre compose this map with the target deck identification $g_2$ to return the centre $g_2^{-1}z_1$ to $z_1$. Its lattice map becomes $A_2H$ and it conjugates the clockwise generator to the inverse generator. At the second centre its lattice map is $H$ with the same reversal. Consequently the target affine vectors are exactly
\[
 w'_1=-A_2H\varepsilon=-\varepsilon+d,\qquad
 w'_2=-H(-\varepsilon')=\varepsilon'-3d.
 \label{I:eq:real-reflected-affine-data}
\]
The minus signs arise because the inverse of a generator with invariant translation $w'_j/m_j$ has translation $-w'_j/m_j$. This calculation verifies the full affine actions on the root covers. For the overlaps choose the root coordinates of the real local base coordinates $t$ and $t-1$ compatible with these local reflections. This choice has an explicit comparison with any original root coordinate: their ratio is a holomorphic unit, so changing it changes the collar by $(2\pi i)^{-1}\log h\,\Pi w_j$, a translation extending on the entire root family. Its inverse gives a holomorphic isomorphism of the two glued descriptions. Thus the extension property and final classification transfer to the original coordinates. Then $\log s'_j=\overline{\log s_j}$ on paired branches. Under $\zeta\mapsto-\bar\zeta$, conjugation of $1/(2\pi i)$ supplies the compensating sign, and the identity of the two logarithmic collar translations follows exactly. Two roots of the same fixed real coordinate differ by an $m_j$th root of unity; this changes the comparison only by a constant period-coordinate translation, and the additive correction of \ref{fam:src:A.3.1} glues those constants. The preceding holomorphic-unit comparison handles a different original local base coordinate. Its inverse has the same construction.

\item At the cusp choose $s=(2\pi i)^{-1}\log t_c$, so $s\mapsto-\bar s$. The full holomorphic unit matrix satisfies
\[
 C_{-\bar c}(\bar t_c)=-\overline{C_c(t_c)}+
       \begin{pmatrix}6&0\\0&1\end{pmatrix}.
\]
The integral extra matrix has trivial exponential. Hence $(x,t_c)\mapsto(\bar x,\bar t_c)$ intertwines the full cusp deck groups, with the same integral index $\lambda$. The toric fan is defined over $\mathbb Z$, so this map extends antiholomorphically on every toric chart and stratum. It preserves the untwisted cusp collar. We have therefore constructed an antiholomorphic map to the target datum with parameter $-\bar c$, $r'_1=-1,k'_1=1,r'_2=1,k'_2=-3,l'_0=0$, whose $p'$ is $+1$.

\item Apply \ref{fam:src:B.2.1} to normalize this target back to the canonical $p=-1$ datum. Take $a=-1$ and the odd shift $n=1$. This gives a holomorphic map to $X_{1-\bar c}$. Every other antiholomorphic map, composed with the inverse of this one, is a holomorphic comparison between canonical members. \ref{fam:src:A.5.1} gives exactly the stated congruence. Exponentiating it gives $w'=-\bar w$.

\item We construct an involutive splitting of the group extension. For a real member, choose any resulting antiholomorphic self-map $J$. Its derivative on ordinary fibre coordinates is $+\text{complex conjugation}$: the initial reflection has derivative $-\text{complex conjugation}$, while the normalization just used has linear part $-I_2$. Hence it conjugates translation by $b\in\mathbb C/\mathbb Z$ to translation by $\bar b$. Its square is a holomorphic translation $b_0$. Commuting $J$ with $J^2$ shows $b_0=\bar b_0$ in $\mathbb C/\mathbb Z$. This equality forces the imaginary part of a representative of $b_0$ to vanish, so take $b_0\in\mathbb R$. Replacing $J$ by the composition of translation by $-b_0/2$ with $J$ changes its square to $b_0-b_0/2-b_0/2=0$. We have an antiholomorphic involution. Under $b\mapsto e^{2\pi ib}$, its conjugation action is $z\mapsto1/\bar z$. This proves the split group and specifies its antiholomorphic coset.
\end{enumerate}
\end{proof}

\begin{proposition}[Two conjugacy classes of real structures]\hypertarget{tag-02PE}{}\label{mapsI:src:A.3.1}
\StatementRecord{mapsI:src:A.3.1}{proved}{02PE}{\SeriesRef{fam:src:B.4.2}, \SeriesRef{fam:src:B.4.1}; \cite[Theorem~3.1, p.~48]{Illman1972}; \cite[Theorem 4.1(1)]{Knudson2004}}

For a real member the antiholomorphic involutions have exactly two
conjugacy classes under holomorphic automorphisms. For any fixed
involutive splitting $\sigma$ they are represented by $\sigma$ and
$\rho(1/2)\sigma$. Every nonempty fixed set is a closed connected
orientable totally real three-manifold with the mod-two homology of
$S^3$. The circle $\rho$ preserves each such fixed set and acts
semifreely on it.
\end{proposition}



\begin{proof}
\begin{enumerate}[label=\arabic*.,leftmargin=*]
\item Proposition~\ref{fam:src:B.4.2} identifies the full holomorphic group
with $\mathbb C^*$ and gives the conjugation action $z\mapsto1/\bar z$.
Thus the square of $z\sigma$ is $z/\bar z$, and it is an involution
exactly for $z\in\mathbb R^*$. Conjugation by $y\in\mathbb C^*$ changes
this coefficient to $|y|^2z$, giving exactly the two sign classes.

\item At a fixed point an antiholomorphic involution has a real three-dimensional
fixed tangent space, and its complexification is the full tangent space.
An invariant metric and its exponential map identify its fixed locus
locally with this real subspace. It is therefore a closed totally real
three-manifold. The smooth $\mathbb Z/2$ action admits an equivariant
triangulation by \cite[Theorem~3.1, p.~48]{Illman1972}; compactness and finite
dimension give the finite-dimensional $G$-CW hypothesis. Smith's theorem,
in the form of
\cite[Theorem 4.1(1)]{Knudson2004}, applies to the smooth finite action on
the sphere and gives the stated mod-two homology. A nonempty
three-dimensional fixed set is consequently connected, and its first
mod-two cohomology vanishes; its orientation class $w_1$ therefore
vanishes. On the unit circle $1/\bar z=z$, so $\sigma$ commutes with
$\rho$. Its stabilizers are inherited from the semifree action on $X$ and are
therefore either trivial or the whole circle.
\end{enumerate}
\end{proof}

\begin{definition}[Representatives of the real involutions]\hypertarget{tag-IP05}{}\label{I:def:real-representatives}
Put
\[
 L_F=(2\pi i)^{-1}\log F,\qquad L_G=(2\pi i)^{-1}\log G,
 \qquad F^3=t,\quad G^4=t-1.
 \hypertarget{tag-02P4}{}\label{mapsI:real-eq-1}
\]
The roots are positive over $t>1$. We continue from that interval
through the upper half $t$-plane. Their logarithms are taken on the
orbifold uniformizing cover. Write $\iota$ for its reflection lift.
Every point of $-i\R_{>0}$ has a representative $c=1/2+iy$, and every
point of $i\R_{>0}$ has a representative $c=-1/2+iy$, with $y$ in the
admissible range. Set
\[
\begin{array}{c|c|c}
 &\alpha&S(t)\\\hline
 c=1/2+iy&1/3&2L_F-4L_G-1/6\\
 c=-1/2+iy&1/2&-2L_G+1/4
\end{array}
\qquad b_c(z)=\alpha\Pi_c(z)g+S(t(z))e.
\label{mapsI:real-eq-2}
\]
\end{definition}

\begin{theorem}[Explicit real involutions and all fixed fibres]\hypertarget{tag-02PF}{}\label{mapsI:src:A.3.2}
\StatementRecord{mapsI:src:A.3.2}{proved}{02PF}{\SeriesRef{fam:src:B.4.2}, \SeriesRef{mapsI:real-eq-4}, \SeriesRef{mapsI:real-eq-2}, \SeriesRef{fam:src:A.3.1}, \SeriesRef{mapsI:real-eq-1}}

On the ordinary period charts the formulas
\[
 \sigma_c[z,\zeta]=[\iota z,\bar\zeta+b_c(\iota z)],
 \qquad \sigma_c^- =\rho(1/2)\sigma_c
 \label{mapsI:real-eq-3}
\]
extend to antiholomorphic involutions of the compact threefold.
With these representatives their complete fixed-set data are as follows.
Here $2T^2$ means the disjoint union of two real two-tori; every displayed
circle in a double curve avoids its two triple-point endpoints.
\begin{center}\small
\begin{tabular}{cclccc}
\toprule
$w$&involution&image in $\RP^1$&on $S_1$&on $S_2$&on $W$\\
\midrule
$-ir$, $r>0$&$\sigma_c$&all $\RP^1$&$2T^2$&$2T^2$&$T^2\sqcup S^1_{D_2}\sqcup S^1_{D_{12}}$\\
$-ir$, $r>0$&$\sigma_c^-$&$[1,\infty]$&$\varnothing$&$T^2$&$S^1_{D_1}\sqcup S^1_{D_2}$\\
$ir$, $r>0$&$\sigma_c$&$\varnothing$&$\varnothing$&$\varnothing$&$\varnothing$\\
$ir$, $r>0$&$\sigma_c^-$&$[\infty,1]$ through $0$&$2T^2$&$T^2$&$S^1_{D_1}\sqcup S^1_{D_{12}}$\\
\bottomrule
\end{tabular}
\end{center}
The first interval is the closure of $t>1$; the second is the closure
of $t<1$. On every open real arc contained in the indicated image the
real regular fibre is $2T^2$. No triple point is real for these structures.
\end{theorem}



\begin{proof}
\begin{enumerate}[label=\arabic*.,leftmargin=*]
\item Start with the reflection $(z,\zeta)\mapsto(\iota z,-\bar\zeta)$
and the integral matrix $H$ of \eqref{I:eq:real-reflection-matrix}.
Proposition~\ref{fam:src:B.4.2} proves its period compatibility and gives
the target affine datum \eqref{I:eq:real-reflected-affine-data}.
The normalization from this
target at parameter $c_0=-\bar c$ to the canonical datum at $c_0+1$
has linear part $-I$ and translation
\[
 b^{\rm norm}_{c_0}=\tfrac13\Pi_{c_0}g+(2L_F-4L_G)e.
 \hypertarget{tag-02P5}{}\label{mapsI:real-eq-4}
\]
Its period defects at the two finite generators are
$(0,-2,2,0)^t$ and $(0,0,2,0)^t$. To check them, the logarithmic increments
are $-2/3$ and $1$, and one multiplies the two displayed matrices $A_j$
by $g/3$. The finite difference vectors in the complex translation space are $2e$ and $-4e$, so subtracting
their logarithmic collars leaves holomorphic translations on both full
root discs. At the cusp
$2L_F-4L_G=s/3-(2\pi i)^{-1}\log(1-q)$ and
$\Pi g=(6\mu,-s+C_{21})^t$. The logarithmic terms cancel and give the
holomorphic toric unit
$(2\mu,C_{21}/3-(2\pi i)^{-1}\log(1-q))$.
Thus \eqref{mapsI:real-eq-4} extends over all three fillings, with its inverse.

\item For $c=1/2+iy$ compose this normalization with the initial reflection
and then with the constant translation $e/6$. This gives exactly the
first line of \eqref{mapsI:real-eq-2}. For $c=-1/2+iy$, additionally close the parameter
shift by the following explicit even-shift formula with $h=-2$,
\[
 b_{h,c_0}=-\frac h{12}\Pi_{c_0}g+
 \left(h(L_F-L_G)-\frac{h^2}{24}\right)e.
\]
A final constant translation $-e/4$ gives the second line of \eqref{mapsI:real-eq-2}.
The displayed translation represents the even shift of Proposition~\ref{fam:src:A.3.1}. Its finite-chart corrections are holomorphic. For its cusp expression, put $\ell=L_F-L_G=-s/12+\kappa_0(q)$ and
$C_{c,21}(q)=c+C_{0,21}(q)$, and set
$D_0(q)=C_{0,21}(q)-12\kappa_0(q)$.
The explicit even-shift translation becomes
\[
 b_{h,c}=
 \left(-h\mu/2,-h(c+D_0)/12-h^2/24\right)^t.
\]
The logarithmic terms have cancelled exactly. Exponentiating both
coordinates gives holomorphic nonzero torus multipliers on the whole cusp
disc; multiplication by these units extends over every toric stratum.
Compatibility with the period deck group is the integral identity
$\exp(2\pi i h\lambda_1)=1$ for $h\in2\mathbb Z$ and
$\lambda_1\in\mathbb Z$. The finite-chart defects are checked by the
same displayed $A_j$ and root logarithms used for the normalization.

\item Over $t>1$ we have
$\Re\mu=1/2$, $\Re\tau=\Re\beta^{\rm part}=0$ and the logarithms in
\eqref{mapsI:real-eq-1} are purely imaginary. The two squares are respectively the
integral period translations $2\widehat w$ and $3\widehat w$.
The square is a holomorphic automorphism and is consequently the
identity everywhere. This proves the involution claim.

\item We next give exact equations for all real fibres. Put $a=\Re c$ and
$R=-H+2aE$. In the chosen branch the reflection identifications on
$t>1$, $0<t<1$, $t<0$ are respectively $1,g_2,g_0^{-1}$.
They follow by reflecting the chosen upper-half-plane paths; equivalently
the root arguments are $(0,0)$, $(0,\pi/4)$ and $(\pi/3,\pi/4)$.
The three fibre involution matrices are
\[
 R_>=R,\qquad R_m=A_2^{-1}R,\qquad R_<=M_0R.
 \hypertarget{tag-02P6}{}\label{mapsI:real-eq-5}
\]
In marked real coordinates let $r_c(t)=\Pi_c(t)^{-1}b_c(t)$,
where $\Pi_c^{-1}$ is the real inverse. The fixed tori are exactly
\[
 \{x\in\R^4/\Lambda:(I-R_\bullet)x=
                       r_c(t)+\epsilon d/2\pmod\Lambda\},
 \qquad\epsilon=0,1.
 \label{mapsI:real-eq-6}
\]
This specifies the cosets, including their parameter dependence. Their
identity subgroups are $T_{\langle w,d\rangle}$ on both outer arcs and
$T_{\langle u+w,d\rangle}$ on the middle arc. When nonempty there are
two cosets. Relative to any one solution their offsets are
$(g+u)/2$ on the right arc and $g/2$ on the other two arcs.

\item Existence of fixed points is decided by integral affine congruences. An affine involution
$x\mapsto R_\bullet x+\beta$ has a fixed point if and only if its
square translation $(I+R_\bullet)\beta$ belongs to
$(I+R_\bullet)\Lambda$. In \eqref{mapsI:real-eq-2} the square translation is
\[
 (I+R_\bullet)\alpha g+2\Re S\,d.
\]
For $c=1/2+iy$ the values $2\Re S$ on the three arcs are
$-1/3,-4/3,-2/3$; for $c=-1/2+iy$ they are $1/2,0,0$.
The resulting integral square vectors are respectively
$(2w,0,2w)$ and $(3w,d,3w)$. Adding $d/2$ to the involution adds $d$
to each square vector. Taking column Hermite normal forms of
$I+R_\bullet$ gives exactly the nonempty entries in the table.
All three $I-R_\bullet$ have nonzero Smith factors $1,2$.
This proves the coset counts and offsets. The Hermite and Smith
computations are listed in \cref{tab:I:programs}. The left-arc lift
$g_0^{-1}$ follows from the reflected upper-half-plane path above.

\item At $t=0$ the root degree is odd. Conjugation by the degree-three deck
group acts transitively on the three reflection lifts; the centralizer
of a reflection in that deck group is trivial.
Thus the quotient has exactly the same two real tori as the adjacent
nonempty real regular fibres. In the empty cases there are none. This
also follows by solving the limiting affine equations on the full
multiplicity-three root quotient.

\item At $t=1$ write the affine deck map as
$G(x)=A_2x+v_2/4$. The real root involution has matrix $R_>$ and
translation
\[
 b_-=(1/3,0,0,-1/6)^t,
 \qquad b_+=(1/2,0,0,1/4)^t
\]
for the two parameter signs. The real locus of the quotient is the
quotient of the fixed sets of $G^kh$, $0\leq k<4$. Its even and odd
lift parities are identified internally by $G$, while $G^2$ preserves
each parity. Direct integer congruences show that the nonempty lift
parities are: both for negative imaginary $w$ and $\epsilon=0$; only even for
negative imaginary $w$ and $\epsilon=1$; neither for positive imaginary $w$ and $\epsilon=0$;
only odd for positive imaginary $w$ and $\epsilon=1$. Every nonempty lift has
two tori, exchanged by $G^2$. Each parity therefore contributes one
torus downstairs. The congruence enumeration uses denominator $24$:
translation denominators divide $12$, and the nonzero Smith factors are
$1,2$, so every component has such a representative. The component
permutation is included in the exact enumeration listed in
\cref{tab:I:programs}. The normal
real line has fourth power of sign positive for the even parity and
negative for the odd parity. This establishes the two one-sided
endpoints at $t=1$ recorded in the table.

\item We compute the cusp fixed set. Put
$A=C_{12}(0)$; the cusp normalization gives
\[
 \Re\mu(0)=1/2,\quad A\in i\R,\quad C_{21}(0)=c-A,
 \quad C_{11}(0)=6\mu(0),\quad C_{22}(0)=\mu(0).
\]
In a toric lift the involution is
$x\mapsto U/\bar x$, with
\[
\begin{array}{c|cc}
 &b_1(0)&b_2(0)\\\hline
 w=-ir&2\mu(0)&C_{21}(0)/3-1/6\\
 w=ir&3\mu(0)&C_{21}(0)/2+1/4
\end{array}
\qquad U_j=e^{2\pi i b_j(0)}.
\hypertarget{tag-02P7}{}\label{mapsI:real-eq-7}
\]
Thus $U_2>0$ in both rows, whereas $U_1>0$ in the first and $U_1<0$
in the second. Replacing $\sigma_c$ by $\sigma_c^-$ negates $U_2$.
The smooth central orbit has a real torus precisely when both units
are positive, with equations $|x_j|^2=U_j$.

\item For an edge $v$, set $n=(-v_2,v_1)^t$. The period translate bringing
the inverted edge back is $B_0n=v$. Its curve coordinate $z=x^n$ has
involution
\[
 z\longmapsto K_v/\bar z,
 \qquad K_v=U^n e^{2\pi i n^tC(0)n}.
 \hypertarget{tag-02P8}{}\label{mapsI:real-eq-8}
\]
The signs for $(D_1,D_2,D_{12})$ are $(-,+,+)$ for negative $w$
and $(-,-,-)$ for positive $w$. Negating $U_2$ reverses the first and
third signs, but not the second. A positive sign gives the single
circle $|z|^2=K_v$; a negative sign gives no fixed point. Central
inversion exchanges the two triangle classes of the periodic fan, so
it exchanges the two triple points. These calculations prove every
cusp entry. Finally a nonempty fixed three-manifold cannot be confined
to the special fibres: a holomorphic function constant on an open
piece of a totally real submanifold is locally constant. This also
proves that the positive-$w$ representative with no regular real
fibres is entirely fixed-point free.
\end{enumerate}
\end{proof}

\begin{theorem}[The real three-manifolds]\hypertarget{tag-02PG}{}\label{mapsI:src:A.3.3}
\StatementRecord{mapsI:src:A.3.3}{proved}{02PG}{\SeriesRef{mapsI:src:A.3.2}}

Every nonempty fixed set in Theorem~\ref{mapsI:src:A.3.2} is diffeomorphic to $S^3$.
Its fundamental group is trivial and its integral homology is
$(\Z,0,0,\Z)$. For negative imaginary $w$, the restricted circle action
is the standard semifree action on $S^3$ with fixed circle
$D_2^{\sigma_c}$. For positive imaginary $w$, the nonempty real form
has a free circle action isomorphic, up to circle and orientation
conventions, to the Hopf action.
\end{theorem}



\begin{proof}
\begin{enumerate}[label=\arabic*.,leftmargin=*]
\item For negative imaginary $w$ both fixed manifolds meet $D_2$ in one circle.
Their circle action is semifree, so the orbit space is a compact
orientable surface with boundary. The local quotient at a fixed circle
is the quotient of $S^1\times\C$ by scalar rotation on $\C$, hence is
a boundary collar. Over the surface interior the principal circle bundle
is trivial because the surface has boundary and $H^2=0$. Changes of
boundary trivialization extend over the attached normal discs by
rotation. Thus the three-manifold is
\[
 (B\times S^1)\ \cup_{\partial B\times S^1}\ (\partial B\times D^2).
\]
If $B$ has genus $g_B$ and $b_B$ boundary components, van Kampen or its
surface handle decomposition gives first homology
$\Z^{2g_B+b_B-1}$. Smith's theorem makes this rank zero. Since $b_B>0$,
we have $g_B=0,b_B=1$. The displayed manifold is then
$\partial(D^2\times D^2)=S^3$, with the indicated circle action. This
derives the smooth recognition from an explicit surface decomposition.

\item For positive imaginary $w$ the nonempty real manifold is free under
$\rho$. The two real tori on $t<1$ join at the single real torus over
$1$. At $0$ the odd root comparison is a torus product. At infinity
they are capped by the real neighbourhoods of $D_1$ and $D_{12}$.
These are solid tori: the primitive circle in the edge direction
collapses and the remaining real circle is their core. In the outer
real-fibre lattice $\Z\langle w,d\rangle$ the two cap meridians are
$w$ and $w-d$.

\item We check the intervening torus identification. Put
$B_1=A_2^{-1}A_1^{-1}A_2$ and $C=B_1^2$. The two real-arc involutions
satisfy $R_<R_m=B_1$, $R_mB_1R_m=B_1^{-1}$ and $B_1^3=I$.
Moving between the real root lifts requires
$B_1^kR_mB_1^{-k}=B_1R_m$, so $2k\equiv1\pmod3$ and $k=2$.
Thus the odd-root transport is $C$. It satisfies $R_<=CR_mC^{-1}$, and
$C(u+w)=-w+d$. The exchange across the multiplicity-four real torus
is $A_2^2$; after this odd-root transport it is
$CA_2^2C^{-1}$, whose restriction to $\langle w,d\rangle$ is
$w\mapsto-w$, $d\mapsto d$. Consequently the meridians, compared on
one boundary torus, have determinant of absolute value one. The
resulting genus-one Heegaard splitting is $S^3$: primitive meridians
intersecting once can be taken to the two coordinate circles by an
integral torus automorphism, giving the usual two solid tori of $S^3$.
The free circle $d$ then has Euler number of absolute value one, so
the displayed bundle over $S^2$ is the Hopf bundle. This proves the
fundamental-group and integral-homology assertions.
\end{enumerate}
\end{proof}

\begin{proposition}[Restrictions of the circle quotient]\hypertarget{tag-02PH}{}\label{mapsI:src:A.3.4}
\StatementRecord{mapsI:src:A.3.4}{proved}{02PH}{\SeriesRef{mapsI:src:A.2.1}, \SeriesRef{mapsI:real-eq-2}}

The map $e_{\rm q}$ of Theorem~\ref{mapsI:src:A.2.1} has a $\rho$-invariant
cusp replacement. Its restriction to
every nonempty real form is null-homotopic. On an ordinary real fibre
its image consists of the two quotient-circle values
\[
 \theta=\tfrac12\gamma(r_c(t)),\qquad
 \theta=\tfrac12\gamma(r_c(t))+\tfrac12\pmod\Z,
 \quad
 \gamma(r_c(t))=\alpha+\frac{\Im S(t)}{D_c(t)},
 \label{mapsI:real-eq-9}
\]
with $\alpha,S$ from \eqref{mapsI:real-eq-2}. The circle action is along each real
two-torus and is annihilated by $e_{\rm q}$.
\end{proposition}



\begin{proof}
\begin{enumerate}[label=\arabic*.,leftmargin=*]
\item The literal quotient is circle invariant. In a cusp extension
$\theta':N_0\to S^1$, each orbit circle has degree zero: its class is
$d$, which vanishes in $\pi_1(N_0)$. The quotients
$\theta'(\rho(s)x)/\theta'(x)$ therefore have a continuous real lift
vanishing at $s=0$ and periodic in $s$. Averaging this real lift over
$s\in\R/\Z$ changes $\theta'$ to a circle-invariant map. The cocycle
identity proves invariance, and on an invariant outer collar the lift
is zero, so the replacement agrees exactly with the prescribed map.
The same construction is smooth. Thus $e_{\rm q}|_{\operatorname{Fix}\sigma}$
factors through its circle orbit space, which is $D^2$ or $S^2$ by the
preceding proof. Both have only two-dimensional homotopy type and
$S^3$ is two-connected, so this restriction is null-homotopic.
The first row of $I-R_\bullet$ is $(2,0,0,0)$, giving the two values
in \eqref{mapsI:real-eq-9}. The real inverse of the period matrix sends $S e$ to
first coordinate $\Im S/D_c$; this follows by subtracting
$(\Im\mu/\Im\tau)$ times its first imaginary-coordinate equation
from the second. Both real-fibre identity subgroups have zero
$\gamma$-coordinate, proving the remaining assertions.
\end{enumerate}
\end{proof}




